% sec:spec-convergence — Spectral Convergence
\section{Spectral Convergence}
\label{sec:spec-convergence}

The previous sections built the machinery of spectral methods ---
polynomial bases, differentiation matrices, quadrature rules --- but
left a central claim unsubstantiated: that spectral methods converge
exponentially for smooth functions.  Finite-difference methods gain a
fixed number of digits of accuracy each time the grid is refined by a
constant factor; spectral methods gain a fixed number of digits each
time a single node is added.  This qualitative difference --- algebraic
versus exponential convergence --- is the reason that twenty CGL nodes
can match the accuracy of hundreds of finite-difference grid points
(\cref{fig:diffmat-convergence}).  The mechanism is the rapid decay of
Chebyshev expansion coefficients, and that decay is governed entirely
by the smoothness and analyticity of the function being approximated.

\paragraph{Coefficient decay for smooth functions.}

The Chebyshev expansion coefficients $a_n$ defined in
\cref{eq:chebyshev-expansion} are Fourier cosine coefficients in the
variable $\theta = \arccos x$.  Integration by parts applied to the
integral $a_n = (2/\pi c_n)\int_0^\pi
f(\cos\theta)\cos(n\theta)\,\d\theta$ (with $c_0 = 2$, $c_n = 1$ for
$n \geq 1$) yields a factor of $1/n$ for
each derivative of $f$ that exists and is continuous.  A function
$f \in C^k[-1,1]$ with a $(k+1)$-th derivative of bounded variation
therefore has coefficients satisfying
%
\begin{equation}
  \abs{a_n} = \order{n^{-(k+1)}}
  \qquad \text{as } n \to \infty \,.
  \label{eq:coeff-decay-smooth}
\end{equation}
%
Each additional derivative of $f$ buys one more power of decay.  A
$C^\infty$ function has coefficients that decay faster than any
polynomial in $1/n$ --- superalgebraic decay --- because the
integration-by-parts argument can be applied arbitrarily many times.
\physnote{The integration-by-parts argument is the same mechanism that
makes Fourier coefficients of smooth periodic functions decay rapidly.
The Chebyshev case inherits this directly through the
$x = \cos\theta$ substitution.}

\paragraph{Geometric convergence for analytic functions.}

Functions that are merely smooth have superalgebraic coefficient decay,
but analytic functions do better: their coefficients decay
geometrically.  The key idea is that an analytic function can be
extended into the complex plane, and the Chebyshev coefficients can
be estimated using contour integrals on an elliptical contour rather
than on $[-1,1]$ itself.

Under the substitution $x = \cos\theta$, the interval
$[-1,1]$ lifts to the real axis of the $\theta$-plane, and the
substitution $z = \cos\zeta$ (with $\zeta = \theta + i\sigma$)
extends the map into the complex $x$-plane.  The image of the strip
$\abs{\operatorname{Im}\zeta} < \sigma_0$ under $x = \cos\zeta$ is
an ellipse in the complex $x$-plane with foci at $\pm 1$ and
semi-axis sum
%
\begin{equation}
  \rho = e^{\sigma_0} \,.
  \label{eq:bernstein-rho}
\end{equation}
%
This ellipse, called the \emph{Bernstein ellipse} $\mathcal{E}_\rho$,
has semi-major axis $a_\rho = (\rho + \rho^{-1})/2$ along the real
axis and semi-minor axis $b_\rho = (\rho - \rho^{-1})/2$ along the
imaginary axis.  When $\rho$ is close to $1$ the ellipse collapses
onto $[-1,1]$; as $\rho$ grows, the ellipse inflates into the complex
plane.
\histnote{The connection between analyticity in a Bernstein ellipse and
Chebyshev coefficient decay appears in Bernstein's work from the early
twentieth century.  The modern treatment follows Trefethen
\cite{Trefethen2000} and the detailed exposition in
\cite{Trefethen2019}.}%
\warnnote{The symbol $\rho$ here denotes the Bernstein ellipse
parameter, not the convergence ratio of a finite-difference stencil
(\cref{sec:fd-convergence}).  Both usages are standard in their
respective fields.}

\paragraph{The geometric decay bound.}

If $f$ is analytic and bounded inside the Bernstein ellipse
$\mathcal{E}_\rho$ for some $\rho > 1$, then its Chebyshev
coefficients satisfy
%
\begin{equation}
  \abs{a_n} \leq \frac{2M}{\rho^n} \,,
  \label{eq:coeff-decay-analytic}
\end{equation}
%
where $M = \max_{z \in \mathcal{E}_\rho}\abs{f(z)}$.
The proof uses Cauchy's integral formula: the Chebyshev coefficient
$a_n$ is a contour integral of $f(z)T_n(z)$ around $[-1,1]$, and
deforming the contour to $\mathcal{E}_\rho$ introduces the factor
$\rho^{-n}$ from the growth of $T_n$ on the ellipse
($\abs{T_n(z)} \leq \rho^n/2$ on $\mathcal{E}_\rho$ for $n \geq 1$).  The
coefficients decay like powers of $1/\rho$ --- geometric convergence,
with the rate set by the size of the largest ellipse inside which $f$
remains analytic.
\warnnote{The constant $\rho$ is determined by the nearest singularity
of $f$ in the complex plane, not by any property of $f$ on the real
interval alone.  A function that appears perfectly smooth on $[-1,1]$
can have slow spectral convergence if it has a pole or branch point
near the real axis.}

\paragraph{From coefficients to interpolation error.}

The Chebyshev interpolant $p_N(x) = \sum_{n=0}^{N}
\hat{a}_n\,T_n(x)$ reproduces the first $N+1$ coefficients of $f$
exactly (up to aliasing), so the interpolation error is controlled by
the tail of the expansion.  For an analytic function bounded in
$\mathcal{E}_\rho$, the truncation error satisfies
%
\begin{equation}
  \norm{f - p_N}_\infty
    \leq \frac{2M\,\rho^{-N}}{\rho - 1} \,,
  \label{eq:spec-interp-error}
\end{equation}
%
since $\sum_{n=N+1}^{\infty} \abs{a_n} \leq
2M\sum_{n=N+1}^{\infty}\rho^{-n} = 2M\rho^{-N}/(\rho - 1)$ and
$\abs{T_n(x)} \leq 1$ on $[-1,1]$.  The error decreases by a
factor of $\rho$ for each additional node, confirming the exponential
convergence observed in \cref{fig:diffmat-convergence}: the gold curve
drops by roughly a constant factor per node until reaching machine
precision.  The aliasing error from the discrete projection adds terms
of order $\rho^{-2N}$, which is negligible compared to the truncation
tail.
\xrefnote{The Lebesgue constant (\cref{eq:lebesgue-bound}) provides an
alternative route to the interpolation error: for CGL nodes,
$\Lambda_N \sim (2/\pi)\ln N$, so the interpolant is within a
logarithmic factor of the best polynomial approximation.  The bound
\cref{eq:spec-interp-error} is tighter for analytic functions because
it exploits the specific decay rate of the coefficients.}

\paragraph{Runge's function revisited.}

Runge's function $f(x) = 1/(1 + 25x^2)$ from
\cref{sec:spec-interpolation} provides a concrete illustration.  Its
singularities are the simple poles at $x = \pm i/5$, which lie at
distance $1/5$ from the real axis.  The largest Bernstein ellipse inside which $f$ remains analytic has
semi-axis sum $\rho$ determined by $b_\rho = 1/5$, where
$b_\rho = (\rho - \rho^{-1})/2$.  Solving gives
%
\begin{equation}
  \rho = \frac{1}{5} + \sqrt{1 + \frac{1}{25}}
  = \frac{1 + \sqrt{26}}{5}
  \approx 1.220 \,.
  \label{eq:runge-rho}
\end{equation}
%
The poles lie \emph{on} $\mathcal{E}_\rho$, so the bound
\cref{eq:coeff-decay-analytic} holds for any $\rho' < \rho$; in
practice the convergence rate approaches $\rho$ from below.
Each additional Chebyshev node reduces the interpolation error by
roughly the factor $\rho \approx 1.22$, or about $0.86$ digits per
node.
At $N = 15$, the predicted error is roughly $1.22^{-15} \approx
0.037$, which matches the $\approx 4\%$ maximum error visible in
\cref{fig:runge-phenomenon}(b).  The modest convergence rate
--- compared to the machine-precision accuracy achieved at $N = 20$ for
$f(x) = e^{\sin(\pi x)}$ in \cref{fig:diffmat-convergence} ---
reflects the poles lurking close to the real axis.  An entire
function (no finite singularities) has $\rho = \infty$, and the bound
\cref{eq:spec-interp-error} must be replaced by a supergeometric
estimate; for $e^{\sin(\pi x)}$, the coefficients decay
roughly as $\order{(2/(\pi e))^n / \sqrt{n}}$.

\paragraph{The convergence hierarchy.}

The relationship between smoothness and convergence rate is a
hierarchy with sharp boundaries:

\begin{center}
\begin{tabular}{lll}
  \toprule
  Smoothness of $f$ & Coefficient decay & Convergence rate \\
  \midrule
  $C^k$, $(k{+}1)$-th deriv.\ of BV
    & $\abs{a_n} = \order{n^{-(k+1)}}$
    & algebraic, $\order{N^{-(k+1)}}$ \\
  $C^\infty$
    & faster than any $n^{-k}$
    & superalgebraic \\
  Analytic in $\mathcal{E}_\rho$
    & $\abs{a_n} = \order{\rho^{-n}}$
    & geometric, $\order{\rho^{-N}}$ \\
  Entire
    & faster than any $\rho^{-n}$
    & supergeometric \\
  \bottomrule
\end{tabular}
\end{center}

Each step up in smoothness produces qualitatively faster convergence.
A function with a discontinuity in its $k$-th derivative converges at
the algebraic rate $\order{N^{-(k+1)}}$, indistinguishable from a
finite-difference method of order $k + 1$ --- spectral methods buy
nothing for non-smooth functions.  The exponential advantage emerges
only when $f$ is analytic, and the rate improves as the nearest complex
singularity moves further from the real axis.
\physnote{The metric functions of the Schwarzschild and Kerr spacetimes
in Kerr--Schild coordinates are rational functions of the Cartesian
coordinates --- analytic everywhere except at the physical singularity
($r = 0$).  On any domain that excludes the singularity, the spectral
solver achieves geometric convergence with $\rho$ determined by the
distance to $r = 0$ in the mapped coordinate.  This is why the XCTS
solver reaches machine precision at moderate polynomial degrees
($N \sim 20$--$30$).}

\paragraph{Convergence of spectral derivatives.}

Differentiating the Chebyshev expansion term by term introduces
algebraic factors of $n$ in the derivative's coefficients, worsening
the decay by at most one algebraic power.  For an analytic function
with geometric coefficient decay $\abs{a_n} = \order{\rho^{-n}}$,
these algebraic factors are absorbed by the geometric term:
$n\rho^{-n} \to 0$ at the same geometric rate.  The spectral
derivative therefore converges geometrically at essentially the same
rate as the function itself --- the differentiation matrix amplifies
errors by $\order{N^2}$ (\cref{eq:diffmat-norm}), but the
exponentially small interpolation error shrinks fast enough to
overwhelm this amplification.

Quantitatively, the derivative error satisfies
%
\begin{equation}
  \norm{f' - p_N'}_\infty
    \leq C\,N^2\,\rho^{-N}
  \label{eq:spec-deriv-error}
\end{equation}
%
for an analytic function in $\mathcal{E}_\rho$, where $C$ depends on
$\rho$ and $M$ but not on $N$ \cite{Trefethen2000}.  The algebraic
prefactor $N^2$ is negligible compared to the geometric decay for any
$\rho > 1$: at $N = 20$ and $\rho = 2$, the bound is roughly
$400 \times 10^{-6} = 4 \times 10^{-4}$, while the actual error is
typically much smaller because the bound is not tight.  This is the
precise sense in which spectral convergence ``overwhelms'' the
conditioning of the differentiation matrix, as claimed in
\cref{sec:spec-diffmat}.
\warnnote{The geometric convergence of spectral derivatives breaks
down at discontinuities.  If $f$ has a jump in its $k$-th derivative,
the spectral derivative converges only at rate
$\order{N^{-k}}$ --- one power slower than the function itself.
For the DG solver, each element carries its own polynomial, so
discontinuities are confined to element interfaces and handled by
the numerical flux, not by the spectral basis.}
